\begin{textsample}{POS Dim 4 – mg – Score 16.00 – mg/mg\_em\_136.txt}  \label{ex:f4_pos_191}
O mesmo não ocorreu com a Lei 11.645 de 2008, que igualmente altera a LDB nos mesmos artigos.
No entanto, o CNE, em sua manifestação, já antevia, com clareza, que o tema do preconceito, do \textbf{racismo} e da \textbf{discriminação}, se por um lado atinge mais forte e amplamente a \textbf{população} negra, também se volta contra outras formas da diversidade.
Assim, aparece, em diversas passagens, alerta para a necessidade de contemplar a temática \textbf{indígena} em particular, quando se tratar da educação para as relações étnico-raciais.
Referências à Lei 11.645 são feitas, sempre que couber, para que as ações curriculares sejam orientadas ao combate de todas as formas de preconceito, \textbf{racismo} e \textbf{discriminação} que porventura venham a se manifestar no ambiente escolar.
A SEE-MG, à medida que enxerga a educação como um todo, cria as condições necessárias para ampliar a qualidade social do ensino oferecido a nossas crianças, adolescentes, jovens e adultos.
Já foi dito, com razão, que as lutas de libertação libertam também os opressores.
Já foi constatado que as manifestações do preconceito estão amparadas em visões equivocadas de superioridade entre diferentes, transformando diferenças em desigualdades.
Por tudo isso, incluir a temática da Lei 11.645 no Currículo Referência faz justiça às lutas dos movimentos sociais negros no Brasil que desde há muito alertam a sociedade brasileira sobre o reconhecimento do \textbf{racismo} em nossa sociedade e sobre como deve ser combatido firmemente, seja qual for o grupo que sofra a \textbf{discriminação} e o preconceito.
A sociedade brasileira deve ao movimento social negro um tributo por sua coragem em se empenhar, com determinação e persistência, pela construção de uma sociedade nova, onde a diferença seja vista como uma riqueza e não como um pretexto para justificar as desigualdades.
Nos anos de 1930 percebemos as primeiras iniciativas de estabelecimento, no Brasil, de uma educação pluralista e inclusiva.
Verifica -se na Frente Negra Brasileira uma das propostas iniciais para a oferta do conteúdo da História da África e as trajetórias dos povos negros no Brasil a fim de combater práticas discriminatórias sofridas por crianças negras no ambiente escolar.
Abdias do Nascimento, na liderança do Teatro Experimental do Negro ( TEN ), na década de 1940, defendeu a necessidade de uma formação global das pessoas negras como forma de ações afirmativas para minimizar os abismos de desigualdades provocadas pelo período de escravização.
Outros atores tiveram importantes contribuições, como Movimento Negro Unificado ( MNU ) e Movimento Social Negro ( Mulheres Negras Educadoras Negras, os Grupos de Estudos Afro das Faculdades de várias Universidades do Brasil ).
Todos esses movimentos foram essenciais para sistematizar uma proposta que pudesse ser implementada em todo \textbf{território} nacional, contribuindo para a reflexão e para uma nova prática acerca da diversidade étnico-cultural existente de nosso país, afastando essencialismos, \textbf{discriminações} e preconceitos.
A III Conferência Mundial contra o \textbf{Racismo}, a \textbf{Discriminação} \textbf{Racial}, a \textbf{Xenofobia} e as Formas Correlatas de \textbf{Intolerância}, realizada em Durban, África do Sul, em 2001, foi um importante marco na agenda internacional, do compromisso de diversos chefes de \textbf{estado}, na luta contra o \textbf{racismo} e outras \textbf{intolerâncias}.
Tal marco na agenda internacional provoca no Brasil uma ampliação dos debates das relações \textbf{raciais}, de \textbf{discriminação} e violência, o que exigiu formas específicas de desenhar propostas em prol da equidade para a \textbf{população} negra.
Ao criar a Secretaria Especial de Políticas de Promoção da Igualdade \textbf{Racial} ( SEPPIR ), em 2003, a construção de ações de uma educação para as relações \textbf{raciais} é tomada como prioridade na pauta política, o que contribuiu para a criação de leis, programas e estratégias para a efetiva promoção da igualdade \textbf{racial}.
O Parecer Nº CNE/CP 003/2004, de maneira contributiva à implementação da Lei 10.639, levanta pontos importantes que devem ser considerados para os desenvolvimentos dos pressupostos da Lei dentro da dinâmica escolar, ressaltando que todos os conteúdos disciplinares devem contemplar os fundamentos que regem esse marco legal na educação nacional.
Ademais, o Parecer reforça o entendimento de que a Lei trata, em especial, de uma política curricular orientada por uma dimensão histórica, filosófica, social, \textbf{antropológica} da realidade brasileira, além do combate ao \textbf{racismo} e às \textbf{discriminações} que afligem a \textbf{população} negra e \textbf{indígena}.
Logo, um dos objetivos da legislação é que as diferentes entidades de educação desenvolvam dentro de suas práticas pedagógicas uma pedagogia de combate ao \textbf{racismo} e à \textbf{discriminação} elaborada com o objetivo de educação das relações étnico/raciais positiva ( BRASIL, 2013, p. 90 ).
As diretrizes da Resolução CNE/CP 01/2004 consolidam o histórico de luta para o desenvolvimento de uma educação representativa da sociedade brasileira, no enfrentamento aos preconceitos e violências em função da \textbf{raça} e/ou etnia.
Além disso, contribuem para uma mudança paradigmática da produção de conhecimento que outrora privilegiava, majoritariamente, conteúdos e referências europeias, negligenciando as ciências e saberes africanos, afro-brasileiro e \textbf{indígena}.
Logo, as diretrizes, conforme o Art. 2º do CNE/CP 01/2004, são constituídas de orientações, princípios e fundamentos para o planejamento, execução e avaliação da Educação ( BRASIL, 2013, p. 77 ).
Tais diretrizes promovem uma educação cidadã de atuação democrática e consciente da composição multicultural e pluriétnica brasileira na busca de uma educação para as relações étnico-raciais positiva.
É importante destacar que tanto o Parecer Nº CNE/CP 003/2004 quanto a Resolução CNE/CP 01/2004 não trazem em sua redação as questões relacionadas aos povos \textbf{indígenas}.
Tais documentos foram desenvolvidos a partir da Lei 10.639/2003, quando ainda não havia sido promulgada a Lei 11.645/2008, que atualiza a normativa incluindo a perspectiva \textbf{indígena}.
E após essa atualização não foram feitas outras legislações dando tratos específicos das questões \textbf{indígenas} na Educação das Relações Étnico-Raciais.
Todavia, deve -se ter atenção às histórias, culturas, \textbf{identidades} e contribuições \textbf{indígenas} na conformação da nação brasileira a fim de atender os pressupostos da Lei 11.645 e desenvolver — em todo \textbf{território} nacional, nos diferentes níveis e modalidades — uma educação antirracista, multicultural, pluriétnica e atenta aos princípios democráticos de igualdade e equidade.
A Lei 10.639, como dito anteriormente, marca historicamente o início de uma série de ações afirmativas, o que exigiu a implementação de programas complementares a fim de conseguir os pressupostos da legislação.
Foram realizadas pelo Ministério da Educação : formações, presenciais e a distância, de professores na temática de diversidade étnico-raciais ; publicação de materiais didáticos específicos ; realização e fomento à pesquisa do tema ; fortalecimento dos Núcleos de Estudos Afro-Brasileiros ( NEAB ) ; criação de Fóruns Estaduais e Municipais de Educação e Diversidade Étnico-Racial ; publicação específica da lei dentro da Coleção Educação para Todos ; dentre as outras ações.
Alguns programas direcionados ao Ensino Superior intentaram reverter quadro histórico de pouco acesso da \textbf{população} negra nas universidades como, por exemplo, o Programa Diversidade na Universidade ( 2002 ), em colaboração internacional junto à UNESCO, defendeu a inclusão \textbf{racial} e o combate à \textbf{exclusão} social que afastava a \textbf{população} afrodescente e \textbf{indígena} do Ensino Superior, significando na melhoria de condições e as oportunidades para esses grupos.
Em 2012, a Lei 12.711 instituiu reservas de vagas para ingresso nas universidades federais e nos institutos técnicos federais brasileiros, o que significou um importante passo para a diminuição das desigualdades de acesso nestas instituições de ensino.
O \textbf{racismo} afeta profundamente a qualidade das instituições educacionais, prejudicando a trajetória escolar e comprometendo a garantia do direito humano à educação de \textbf{milhões} de crianças, adolescentes, jovens e adultos de nosso país.
Enfrentá -lo é um desafio de toda sociedade brasileira, conforme destaca o Parecer das Diretrizes Curriculares Nacionais para a Educação das Relações Étnico-Raciais e para o Ensino de História e Cultura Afro-Brasileira⁴⁰ e Africana, aprovado pelo Conselho Nacional de Educação em 2004.
Outros programas contribuíram na formação de professores da educação básica para o trato das temáticas étnico-raciais.
O Programa Cultura Afro-Brasileira ( 2005-2006 ) propiciou assistência financeira aos professores e à produção de materiais didáticos específicos ; o Programa UNIAFRO promoveu a formação de nove mil profissionais de educação entre os anos de 2008 e 2010 ; o mesmo programa realizou pesquisas, seminários e publicações acadêmicas a respeito da Lei 10.639, entre os anos de 2005 e 2011 ; a formação a distância por meio dos cursos Educação-Africanidades-Brasil e História da Cultura Afro-Brasileira e Africana contemplou mais de 10 mil profissionais da rede pública, nos anos de 2006 e 2007, dentre as outras ações.
Nesse sentido, verificam -se importantes ações para dar aporte à implementação da lei ao disponibilizar materiais e formações para que essa transformação na educação brasileira se efetivasse.
Em 2015, a Secretaria de Estado de Educação de Minas Gerais cria a campanha “ Afroconsciência ”, dando encaminhamento à implementação da Educação para as Relações Étnico-Raciais no \textbf{Estado} orientada pelas Leis 10.639/03 e 11.645/08, bem como pelo Parecer Nº CNE/CP 003/2004 e pela Resolução CNE/CP 01/2004.
Tal campanha gerou materiais para subsidiar a efetivação das leis nas práticas de ensino da rede estadual de ensino, estimulando diferentes áreas do conhecimento a se inteirar das possibilidades de aplicação dos pressupostos nos diferentes componentes e dos ganhos que esta ação traria para a formação dos estudantes.
Por isso, é importante incentivar e encorajar a equipe pedagógica e docente para que, desde o início dos planejamentos anuais, sejam considerados nos planos de ensino e atividades a presença dos conteúdos de história e cultura africana, afro-brasileira e \textbf{indígena}.
Para além das formações e investimentos realizados para que as Leis 10.639 e 11.645 de fato entrassem e permanecessem nas pedagogias e propostas na Educação Básica e Ensino Superior, é importante que conjuntamente equipe gestora e equipe pedagógica alinhem as propostas para que a Educação das Relações Étnico-Raciais de fato se concretize.
É importante compreender que tais marcos legais buscam que os conhecimentos construídos pelos estudantes, bem como as práticas pedagógicas desenvolvidas na escola, contemplem a diversidade de atores e protagonistas brasileiros, levando em consideração \textbf{raça} e etnia em seus aspectos sociais, culturais e históricos na vida dos cidadãos brasileiros.
Logo, faz -se necessário rever as abordagens e conteúdos que marginalizam os africanos, afrodescentes e \textbf{indígenas} na história de produção de conhecimento no Brasil nas diferentes áreas, desfazendo essencialismos, estereótipos e preconceitos para com seus grupos e suas culturas.
Desse modo, além de enfrentar as \textbf{discriminações} \textbf{raciais} e étnicas e o \textbf{racismo}, a inserção de novas concepções estará estabelecendo uma educação antirracista que valoriza a diversidade e conformação pluriétnica e cultural brasileira. ⁴⁰ O Parecer do Conselho Nacional de Educação ( CNE/CP3/2004 ) foi elaborado pela Professora Doutora Petronilha Beatriz Gonçalves e Silva, professora Emérita da Universidade Federal de São Carlos ( UFSCAR/SP ), intelectual reconhecida nacional e internacionalmente por seu trabalho no campo da educação das relações étnico-raciais.
Por indicação do movimento negro, Petronilha foi a primeira \textbf{mulher} negra a ocupar um assento no Conselho Nacional de Educação ( CNE ), exercendo mandato no período de 2002 a 2006.
Nessa condição, foi relatora do Parecer das Diretrizes Curriculares Nacionais para a Educação das Relações Étnico-Raciais e para o Ensino de História e Cultura Afro-Brasileira e Africana, aprovado pelo CNE, em 2004.

% matched lemmas: antropológico, discriminação, estado, exclusão, identidade, indígena, intolerância, milhão, mulher, população, racial, racismo, raça, território, xenofobia
\end{textsample}
